There are three steps as follows.
\begin{enumerate}
 \item \textbf{Collecting} all \emph{simple transition path}s.
 Every \emph{simple transition path} is loop free and interleaving free. It
 contains only the edges of atomic assignment or guard transitions without interleaving other paths.
 Each of them corresponds to a path in the flatten program in Definition~4.1 in \cite{GulwaniJK09}.
 \item \textbf{Rewriting} the program $c$ by rearranging all \emph{simple transition paths} as the syntax in~\cite{GulwaniJK09} and preserving the semantics.
 \item \textbf{Refining} the program and computing the 
 refined program, $\rprog$ by Algorithm~1 in paper~\cite{GulwaniJK09}.
 This step invokes the algorithm REFINE from paper~\cite{GulwaniJK09} and computes the 
 refined program $\rprog$ for a program $c$ given the rewritten program as input.
\end{enumerate}

\subsection{The Simple Transition Path}
We first collect the loop headers $\loopl(c) \subseteq \progl(c)$ from a program $c$, which is the set of all program points corresponding to the loop headers in program $c$.
\begin{defn}[Loop Headers ($\loopl : \cdom \to \mathcal{P}(\ldom)$)]
 \label{def:loopl}
 \[
 \loopl(c) \triangleq 
 \left\{
 \begin{array}{ll}
 \{\} & {c} = \clabel{\assign x e}^{l} \\
 \loopl({c_1}) \cup \loopl({{c_2}}) & {c} = {c_1};{c_2} \\
 \loopl(c_t) \cup \loopl({{c_f}}) & {c} =\eif(\clabel{\bexpr}^{l}, c_t, c_f) \\
 \loopl(c_w) \cup \{l\}, & {c} = \ewhile \clabel{\bexpr}^{l} \edo (c_w)
 \end{array}
\right.
\]
 \end{defn}

\begin{defn}[Simple Transition Path]
 \label{def:tpath}
A \emph{simple transition path}
$\tpath \in \paths(\absG(c))$ for the program $c$, is either a simple cyclic path, which has the same start- and end-point
or a simple path has either different while loop headers, the program entrance or exit as its start- and end-point
without visiting any loop header inside the path.
\\
Specifically, a path $l_0 \xrightarrow{dc_0} l_1 \xrightarrow{dc_1} \ldots l_n \in \paths(\absG(c))$ with the
vertices sequence $(l_0, \ldots, l_n)$, where $l_i \in \absV(c)$ for every $i = 0, \ldots, n$ and
%
the edge sequence $(e_1, \ldots, e_n)$, where $e_i = (l_{i - 1}, dc_i, l_{i}) \in \absE(c)$ for every $i = 1, \ldots, n$,
%
is a \emph{simple transition path} if and only if it satisfies,
\begin{itemize}
 \item $l_i \neq l_j$ for every $i = 0, \ldots, n$ and $j = 1, \ldots, {n - 1}$,
 \item $l_0$ is either the program point of a loop header or the program entrance ($l_0 = 0$),
 i.e., $l_0 \in \loopl(c) \cup \{ 0 \}$
 \item and $l_n$ is either the program point of a loop header or the program exit ($l_n = \lex$),
 i.e., $l_0 \in \loopl(c) \cup \{ \lex \}$,
 \item and $l_i \notin \loopl(c) \cup \{ 0, \lex \}$ for every $i = 1, \ldots, n-1$.
\end{itemize}
\end{defn}


\paragraph{Example.}
% \todo{The walk through example}
For example in Figure~\ref{fig:relatedNestedWhileOdd-overview}(b), the $1 \to 2 \to 8 \to 1$ is a \emph{simple transition path}.
However, $1 \to 2 \to 3 \to 4 \to 5 \to 4 \to 6 \to 7 \to 1$ is not a \emph{simple transition path} because it is not simple.
Also, $1 \to 2 \to 3 \to 4 \to 6 \to 7 \to 1$ is not a \emph{simple transition path} because it visits a loop header $4$ inside the path. $1 \to 2 \to 3$ isn't either because its end-point isn't a loop header.
In Figure~\ref{fig:threeWhile-overview}(b), $1 \to 2 \to 3 \to 4 \to 5 \to 6$ is not a \emph{simple transition path} on $\absG(\kw{threeNestedWhile}(n, m, N))$ because it visits a loop header $3$ inside the path.

\subsection{The Program Rewriting}
\paragraph{Syntax.}

\[
    \rprog := \tpath ~|~ \rprepeat(\rprog) ~|~ \rprog;\rprog ~|~ \rpchoose{\rprog, \ldots} 
\]

\paragraph{Algorithm.}

\begin{algorithm}
 \caption{Program Rewriting $\kw{Rewrite}$}
 \label{alg:alg-refine_rewrite}
 \begin{algorithmic}[1]
 \REQUIRE program $c$
 \STATE collects all $c$'s \emph{simple transition path}s from $\absG(c)$, $\tpath_1, \ldots, \tpath_n \in \paths(\absG(c))$.
 \STATE \textbf{init}: candidate set $W = \{c_1, \ldots, c_n\}$, where $c_i = \tpath_i$ and $i = 1, \ldots, n$
 \STATE \textbf{while} $W.size()> 1$:
 \STATE \quad create $c' = \rpchoose{c_1, \ldots, c_m}$ 
 s.t. $c_i \in W \land c_i[0] = c_j[0] \land c_i[-1] = c_j[-1], i, j = 1, \ldots, m$.
 \\ \quad $W.add(c')$ \qquad $W.remove(c_1, \ldots, c_m)$
 \STATE
 \quad create $c' = \rprepeat(c)$ s.t. $c_i \in W \land c[0] = c[-1] \land c[0] \in \loopl(c)$
 \\ \quad $W.add(c')$, \qquad $W.remove(c)$
 \STATE \quad create $c' = c_1; c_2$ s.t. $c_1, c_2 \in W \land c_1[-1] = c_2[0]$
 \\
 \quad $W.add(c')$ \qquad $W.remove(c_1, c_2)$
 \STATE \textbf{Endwhile}
 \\ $c^T = W[0]$
 \RETURN $c^T$.
\end{algorithmic}
\end{algorithm}
%
Algorithm~\ref{alg:alg-refine_rewrite} transforms the syntax of the program following~\cite{GulwaniJK09} and preserves the semantics.
We first use a simple depth-first search strategy that collects all the \emph{simple transition path}s satisfying the Definition~\ref{def:tpath}. It guarantees that every $\tpath$ is equivalent to a path $\rho$ in Definition~4.1 of \cite{GulwaniJK09}.
Then
in Line-2, we initialize each candidate $c_i$ with a \emph{simple transition path} $\tpath_i$. New candidates generated in Line-4, 5, and 6 correspond to the if,
while, and sequence statements in paper~\cite{GulwaniJK09} respectively.
\begin{itemize}
 \item
 Line-4: for all the candidates $c_1, \ldots, c_m$ having the same starting and ending vertices, rewrite them into if statement as~\cite{GulwaniJK09}.
 \item
 Line-5: for every candidate $c'$, if it starts and ends with the same vertex, rewrite it into a while loop statement as~\cite{GulwaniJK09}.
 \item
 Line-6: for every two candidates $c_1, c_2$, if $c_1$ ends with the same vertex as $c_2$'s starting label, rewrite them into sequence statement as~\cite{GulwaniJK09}.
\end{itemize}
% \todo{soundness of the program write}

\subsection{The Program Refinement}
We use the algorithm REFINE from paper~\cite{GulwaniJK09} 
to compute the execution orders of
the simple transition paths that come from the same loop.
For example, two simple transition paths $\tpath_1 = (1 \to 2 \to 3 \to 4)$ and 
$\tpath_4 = (1 \to 2 \to 8 \to 1)$ in the same loop Figure~\ref{fig:relatedNestedWhileOdd-overview}(b) have two possible execution orders:
either $\tpath_1$ executes immediately after the execution of $\tpath_4$, or vice verse.
% and compute the 

Then we construct the refined program by explicating the execution orders of simple transition paths.
% refined program $\rprog$ for a rewritten program $c$.
Simple transition paths can iterate on themselves, so an execution order could
contain iterations of simple transition paths, and sequence of iterations.

For example, the simple transition path $\tpath_3 = (4 \to 5 \to 4)$
the loop $L_4$ in Figure~\ref{fig:relatedNestedWhileOdd-overview}(b) has only one possible execution order,
which is the iteration on itself, $\rprepeat(\tpath_3)$.
Each loop in a refined program contains several execution orders over its original simple transition paths.
For example, the loop $L_1$ in Figure~\ref{fig:relatedNestedWhileOdd-overview}(b)
has two execution orders. In the first one,
$\rprepeat(\tpath_4; \tpath_1; 4:\rprepeat(\tpath_3); \tpath_2)$ 
$\tpath_1$ will execute right after the execution of $\tpath_4$.
Then following the execution of $\tpath_1$, loop $L_4$ finishes full iterations and followed by the execution of $\tpath_2$.
In the second execution order,
$\rprepeat(\tpath_1; 4:\rprepeat(\tpath_3); \tpath_2; \tpath_4)$,
$\tpath_4$ is executed after the execution of $\tpath_1; 4:\rprepeat(\tpath_3); \tpath_2$.
Iteration of one execution order is equivalent to one possible full iterations of the original loop program.
So the reconstructed loop for $L_1$ is 
\[
\rpchoose{ 
 \begin{array}{l}
 1: \rprepeat(\tpath_1; 4:\rprepeat(\tpath_3); \tpath_2; \tpath_4), \\
 1: \rprepeat(\tpath_4; \tpath_1; 4:\rprepeat(\tpath_3); \tpath_2) 
 \end{array}
 }.
\]

\paragraph{Running Example.}
For the walk-through an example in Figure~\ref{fig:relatedNestedWhileOdd-overview}(b),
We first collect all its simple transition paths from $\absG(\kw{relatedNestedWhileOdd}(n, m))$ as
$\tpath_0 = (0 \to 1)$,
$\tpath_1 = (1 \to 2 \to 3 \to 4)$,
$\tpath_2 = (4 \to 6 \to 7 \to 1)$,
$\tpath_3 = (4 \to 5 \to 4)$, $\tpath_4 = (1 \to 2 \to 8 \to 1)$, and $\tpath_5 = (1 \to \lex)$.
Using these simple transition paths, we first rewrite this program by Algorithm~\ref{alg:alg-refine_rewrite} and get the rewritten program
%  corresponding to the syntax in~\cite{GulwaniJK09} as
$ \tpath_0 ; \rpchoose{ 1: \rprepeat(\tpath_1; 4:\rprepeat(\tpath_3); \tpath_2), 
1: \rprepeat(\tpath_4) }; \tpath_5$.
Then we construct the refined program as,
\[
 \tpath_0 ; \rpchoose{ 
 \begin{array}{l}
 1: \rprepeat(\tpath_1; 4:\rprepeat(\tpath_3); \tpath_2; \tpath_4), \\
 1: \rprepeat(\tpath_4; \tpath_1; 4:\rprepeat(\tpath_3); \tpath_2) 
 \end{array}
 }; \tpath_5.
\]

\paragraph{Example II.}
 We present again the example from Figure~\ref{fig:whileTwoCounters-overview}(a) which
%   another example from the .Net base-class library 
to show the necessity of the \emph{path refinement}. 
%   It 
It has a two paths loop
 with different \emph{reachability-bound}s on the locations of different paths.
 %
 %
 \input{chapters/reachability/examples/whileTwoCounters-overview}
 %
%  To know the bounds for locations on different branches of a loop, 
%  it is necessary to know the alternating patterns of the two paths.
 Over its abstract transition graph generated by Section~\ref{sec:psRB-progabs} as Figure~\ref{fig:whileTwoCounters-overview}(b), we collect all its simple transition path as in Figure~\ref{fig:whileTwoCounters-overview}(c).
 Then we transform the loop by explicitly computing the interleaving between paths and
 generates a refined program $\rprog$ as the bottom part of Figure~\ref{fig:whileTwoCounters-overview}.

 We first compute that $\tpath_1$ can iterate on itself, so there is one possible execution sub-order,
 $\rprepeat(\tpath_1)$.
 From this refined program, we can explicitly tell two possible path interleaving patterns.
 In the first one, $\tpath_2$ will be executed after all the iterations of $\tpath_1$ are done, and in the second one,
 the entire loop is done after $\tpath_1$ finishes the iteration.
 Then the following computation based on this refined program can compute more accurate reachability-bounds.
